\begin{thebibliography}{99}

\bibitem{NCCubicA}
\emph{滑らかな三次超曲面のPoisson構造：正規形・退化軌跡・代数的量子化},
\texttt{nc-cubic} 論文A、未公表研究稿 A-v61、2026年11月14日。
\href{https://github.com/uedakazushi/nc-cubic/blob/eabe601c1341777742964584dde5314c0bcf64a9/papers/paper_a.tex}{固定 commit の主原稿}。
支持節は \texttt{papers/support/papers/sections/poisson\_geometry.tex}
の命題 \texttt{prop:pg-represent}, \texttt{prop:pg-jacobi-extension}、
定理 \texttt{thm:pg-normal}、系 \texttt{cor:pg-generic}、および
\texttt{general\_dimension\_universal\_quantization.tex}
の定理 \texttt{thm:pg-all-dimensions}、命題 \texttt{prop:pg-ad-relative-chart}。

\bibitem{NFNCCubicB}
\emph{中心三次非可換変形の一般点における可換復元},
\texttt{nc-cubic} 論文B、未公表研究稿 B-v59、2026年11月14日。
\href{https://github.com/uedakazushi/nc-cubic/blob/eabe601c1341777742964584dde5314c0bcf64a9/papers/paper_b.tex}{固定 commit の主原稿}。
\texttt{papers/generic\_recovery/} の
\texttt{main\_result.tex}, \texttt{generic\_descent.tex},
\texttt{coherent\_saturation.tex}, \texttt{field\_model.tex} を参照。

\bibitem{Threefold}
\emph{非可換三次三様体のAbel圏モジュライ：二次包絡とKuznetsov残余射},
未公表研究稿 v22、2026年11月20日。
\href{https://github.com/uedakazushi/nc-cubic-threefold-moduli/blob/6c35d2c84671c2538e61836a6bf3a511392eacce/papers/nc-cubic-threefold-moduli.tex}{固定 commit の主原稿}。
点スキームは \texttt{papers/sections/point\_scheme\_poisson.tex}、
全 Hom の比較は \texttt{jacobian\_abel\_comparison.tex} を参照。

\bibitem{NFMEG}
P.~Etingof and V.~Ginzburg,
\emph{Noncommutative del Pezzo surfaces and Calabi--Yau algebras},
J. Eur. Math. Soc. \textbf{12} (2010), 1371--1416.
\url{https://doi.org/10.4171/JEMS/235}.
\url{https://ems.press/content/serial-article-files/31746},
Theorem~3.3.2, Section~9.

\bibitem{MoriSmith}
I.~Mori and S.~P.~Smith,
\emph{The classification of 3-Calabi--Yau algebras with 3 generators and 3 quadratic relations},
Math. Z. \textbf{287} (2017), 215--241.
\url{https://doi.org/10.1007/s00209-016-1824-5}.
\url{https://arxiv.org/abs/1502.07403v2}, Theorem~4.2.

\bibitem{LBSV}
L.~Le Bruyn, S.~P.~Smith and M.~Van den Bergh,
\emph{Central extensions of three dimensional Artin--Schelter regular algebras},
Math. Z. \textbf{222} (1996), 171--212.
\url{https://doi.org/10.1007/PL00004532}.
\url{https://sites.math.washington.edu/~smith/Research/Central-extns.pdf},
Theorem~4.2.2, Proposition~4.2.5, Section~4.3.

\bibitem{Mayanskiy}
E.~Mayanskiy,
\emph{Poisson cohomology of two Fano threefolds},
J. Algebra \textbf{424} (2015), 21--45.
\url{https://doi.org/10.1016/j.jalgebra.2014.08.049}.
\url{https://arxiv.org/abs/1303.5916v1}, Theorem~1, Section~2.2.
Theorem~1 は Loray--Pereira--Touzet の分類に帰属されている。
スキーム構造には Section~2.2 の Schouten 計算を用いる。

\bibitem{BCMStandardMonomial}
N.~C.~Bonala, O.~Clarke and F.~Mohammadi,
\emph{Standard monomial theory and toric degenerations of Richardson varieties in the Grassmannian},
\url{https://arxiv.org/abs/2103.16208v1}, 2021,
Section~2.4, Theorem~2.3.

\bibitem{NFMStacksProjective}
The Stacks Project Authors,
\emph{The Stacks Project}, Section~27.21, Projective bundles, Tag~01OA.
\url{https://stacks.math.columbia.edu/tag/01OA}.

\bibitem{NFMStacksAffineDescent}
The Stacks Project Authors,
\emph{The Stacks Project}, Lemma~35.37.1, Tag~0245.
\url{https://stacks.math.columbia.edu/tag/0245}.

\bibitem{NFLZSerre}
X.~Lin and S.~Zhang,
\emph{Serre algebra, matrix factorization and categorical Torelli theorem for hypersurfaces},
Math. Ann. \textbf{391} (2025), 163--177.
\url{https://doi.org/10.1007/s00208-024-02915-8}.
定理番号は \url{https://arxiv.org/abs/2310.09927v2}
の Theorem~3.5, Proposition~3.7, Theorems~4.1 and~5.1 による。

\bibitem{NFBFKKernels}
M.~Ballard, D.~Favero and L.~Katzarkov,
\emph{A category of kernels for equivariant factorizations and its implications for Hodge theory},
Publ. Math. Inst. Hautes \'Etudes Sci. \textbf{120} (2014), 1--111.
\url{https://doi.org/10.1007/s10240-013-0059-9}.
\url{https://numdam.org/item/10.1007/s10240-013-0059-9.pdf},
Theorem~5.39 の一般式とその証明、pp.~82--84。

\bibitem{NFPirozhkov}
D.~Pirozhkov,
\emph{A categorical Torelli theorem for hypersurfaces},
Bull. London Math. Soc. \textbf{56} (2024), no.~10, 3075--3089.
\url{https://doi.org/10.1112/blms.13117}.
本文の Lemmas~5.1 and~5.2、Theorem~5.3 は
\url{https://arxiv.org/abs/2208.13604v1} による。

\bibitem{NFHuybrechtsRennemo}
D.~Huybrechts and J.~V.~Rennemo,
\emph{Hochschild cohomology versus the Jacobian ring and the Torelli theorem for cubic fourfolds},
Algebraic Geometry \textbf{6} (2019), no.~1, 76--99.
\url{https://doi.org/10.14231/AG-2019-005}.
\url{https://content.algebraicgeometry.nl/2019-1/2019-1-005.pdf},
Proposition~3.5 and Remark~3.8.

\bibitem{NFDAGX}
J.~Lurie,
\emph{Derived Algebraic Geometry X: Formal Moduli Problems}, November 5, 2011.
\url{https://www.math.ias.edu/~lurie/papers/DAG-X.pdf},
Section~2, Theorems~5.1.9 and~5.3.16, Proposition~5.3.21.

\bibitem{NFBKP}
A.~Blanc, L.~Katzarkov and P.~Pandit,
\emph{Generators in formal deformations of categories},
Compos. Math. \textbf{154} (2018), 2055--2089.
\url{https://doi.org/10.1112/S0010437X18007303}.
\url{https://arxiv.org/abs/1705.00655v1}, Sections~3.3 and~4.3.

\bibitem{NFKuznetsovFCY}
A.~Kuznetsov,
\emph{Calabi--Yau and fractional Calabi--Yau categories},
J. reine angew. Math. \textbf{753} (2019), 239--267.
\url{https://doi.org/10.1515/crelle-2017-0004}.
\url{https://arxiv.org/abs/1509.07657v2},
Section~3.3, Corollaries~3.18 and~4.1, Lemma~3.19.

\bibitem{NFBFKOrlov}
M.~Ballard, D.~Favero and L.~Katzarkov,
\emph{Orlov spectra: bounds and gaps},
Invent. Math. \textbf{189} (2012), 359--430.
\url{https://doi.org/10.1007/s00222-011-0367-y}.
定理番号は \url{https://arxiv.org/abs/1012.0864v3} の Theorem~5.15 による。

\bibitem{NFKuznetsovCubicFourfolds}
A.~Kuznetsov,
\emph{Derived categories of cubic fourfolds},
in \emph{Cohomological and Geometric Approaches to Rationality Problems},
Progress in Mathematics \textbf{282}, Birkh\"auser, 2010, 219--243.
\url{https://arxiv.org/abs/0808.3351}, Theorem~3.1.
著者稿 \url{https://cims.nyu.edu/~tschinke/books/rationality/kuznetsov.pdf}.

\bibitem{NFMGabrielQuotient}
P.~Gabriel, \emph{Des cat\'egories ab\'eliennes},
Bull. Soc. Math. France \textbf{90} (1962), 323--448,
Chapitre~III, \S1, p.~365.
\url{https://doi.org/10.24033/bsmf.1583}.
\url{https://www.numdam.org/article/BSMF_1962__90__323_0.pdf}.

\bibitem{Fraenkel}
A.~McMillan Fraenkel,
\emph{Extensions of Poisson structures on singular hypersurfaces},
\url{https://arxiv.org/abs/1310.6083v2}, 2013,
Proposition~4.1, Theorem~4.3.

\bibitem{Carles}
R.~Carles,
\emph{D\'eformations dans les sch\'emas d\'efinis par les identit\'es de Jacobi},
Ann. Math. Blaise Pascal \textbf{3} (1996), no.~2, 33--62.
\url{https://doi.org/10.5802/ambp.67}.
\url{https://www.numdam.org/item/AMBP_1996__3_2_33_0.pdf},
Section~8, Lemma~8.35.

\bibitem{CarlesMarquez}
R.~Carles and M.~C.~M\'arquez,
\emph{Different methods for the study of obstructions in the schemes of Jacobi},
Ann. Inst. Fourier \textbf{61} (2011), no.~2, 453--490.
\url{https://doi.org/10.5802/aif.2620}.
\url{https://www.numdam.org/item/AIF_2011__61_2_453_0.pdf},
Section~3, Theorem~3.2, Propositions~3.9 and~3.10.

\end{thebibliography}
